% The six constants of Section 4 on one value axis, and the interval each rank
% closes.  No projection: the drawing is a number line, and the only coordinate
% map is
%   X = 8 v cm       (v the value, so one unit of value is 8cm),
%   Y = the row height in cm, listed below.
% Every abscissa in the body is 8 times the exact constant named in its
% trailing comment, rounded to five places:
%   0                      -> 0.00000
%   1/4                    -> 2.00000
%   1/3                    -> 2.66667
%   1/2                    -> 4.00000
%   h_3 = 3(1-cos 2pi/9)   = 0.7018666706...  -> 5.61493
%   3 - sqrt5              = 0.7639320225...  -> 6.11146
%   alpha_SIC = 2-2/sqrt3  = 0.8452994616...  -> 6.76240
%   (10/9) alpha_SIC       = 0.9392216240...  -> 7.51377
%   1                      -> 8.00000
%   1.2                    -> 9.60000
% The axis carries a minor tick at every tenth of a unit, X = 0, 0.8, ..., 9.6,
% so the scale reads off without a second numeral; the labels are 0, the floors
% 1/4 and 1/3 and 1/2, the three rank-three records h_3 and 3-sqrt5 and
% (10/9)alpha_SIC, the rank-two value alpha_SIC, and the threshold 1.  Labels
% alternate between two rows
% because h_3 and 3-sqrt5 sit 0.4965cm apart at this scale.
%
% EXACT ARITHMETIC.  Each constant is an algebraic number, and four of the six
% are irrational: h_3, 3-sqrt5, alpha_SIC and (10/9)alpha_SIC.  The ledger gives
% a polynomial for three of them, one per independent irrationality.  The fourth
% is omitted because it carries nothing: (10/9)alpha_SIC is a rational multiple
% of alpha_SIC, and its own minimal polynomial 243 z^2 - 1080 z + 800 (least root
% 0.9392216240..., the other 3.5052228204...) is the rescaling of alpha_SIC's.
% Two of the three printed are integral and the third is not; z^2 - 4z + 8/3 is
% quoted in the monic form lem:sicvalue uses, and its integral multiple is
% 3z^2 - 12z + 8:
%   alpha_SIC : z^2 - 4z + 8/3, roots 2 +- 2/sqrt3        (Lemma lem:sicvalue)
%   3 - sqrt5 : s^2 - 6s + 4,   roots 3 +- sqrt5          (Example ex:trine)
%   h_3       : h(s) = 8s^3 - 72s^2 + 162s - 81           (Proposition prop:hesse)
% The ordering drawn is the integer chain of the last ledger line:
%   h(3-sqrt5) = 14 sqrt5 - 27 > 0, since 980 > 729;
%   h(2) = 19 > 0 and the three roots of h sum to 9, so exactly one root lies
%   below 2 and 2 separates the two smallest; 3-sqrt5 < 2 then gives
%   h_3 < 3-sqrt5.
%   12 < 7 sqrt3 gives alpha_SIC > 5/6, hence (10/9) alpha_SIC > 25/27;
%   5 > (223/100)^2 gives 3-sqrt5 < 77/100; and 2500 > 2079 is 25/27 > 77/100.
%   363 < 400 is 11 sqrt3 < 20, which is alpha_SIC < 9/10, i.e. (10/9)alpha_SIC < 1.
% The rational window 7/10 < h_3 < 71/100 of eq:hessewindow is h(7/10) = -136/1000
% and h(71/100) = 588088/10^6.  Every value above was recomputed in exact
% symbolic arithmetic and the paper's printed truncations reproduce to 20
% places.
%
% ROW HEIGHTS, top to bottom; the two real rows are set clear of the four
% complex ones by 0.56 rather than 0.44:
%   3.10  alpha_k(R), k >= 3      2.66  alpha_2(R)
%   2.10  alpha_k(C), k >= 4      1.66  alpha_3(C)
%   1.22  alpha_2(C)              0.78  alpha_1(F)
% A serif of half-height 0.075 ends a closed bracket, an arrowhead an end
% running to 0 with 1/k, a FILLED DOT an end attained by a design, an open dot a
% superseded record.  Five ends carry a filled dot: 1 on the two real rows and on
% alpha_1(F) (thm:corankone(iii)-(iv) at (k+1,k)), alpha_SIC on the alpha_2(C) row
% (lem:sicvalue at (4,2)) and h_3 on the alpha_3(C) row (prop:hesse at (9,3)).  The
% upper end of the alpha_k(C), k >= 4 row carries none, and that is the one real
% distinction the plate draws: (10/9)alpha_SIC is the level thm:ladder proves a
% design to sit BELOW, and no design of that value is exhibited anywhere.  (The
% strictness b < b/(1-w) of the padding remark is a statement about the family of
% bounds, not about which values the padded designs take, and is not the reason.)
% A filled dot is NOT a claim that the row's constant equals that end: whether it
% does on the top row is the conjecture.
% On the alpha_3(C) row, and to the right of its closed end, the
% descent of the rank-three upper bound is drawn as one arrow at the row height
% 1.66 carrying open dots at (10/9)alpha_SIC and at 3-sqrt5 and landing 0.105
% short of h_3: that is Remark rem:sharpc, the record moving down without
% pinning anything.  The two open dots sit under their own guides, so each
% reads against its axis label.
%
% DISCLOSED.  Two marks on the plate are not theorems of this paper.  The
% alpha_2(C) dot is an equality whose LOWER half is Nesterenko's uniform-weight
% bound carried to arbitrary weights in Section 4.4; that half is a citation,
% and it is not formalized -- Gtz.sicDesign_valueAtMost proves "value at most
% alphaRankTwo" and stops short of the attainment, and no attainment
% declaration for the SIC exists (contrast trineMargin_isGreatest and
% hesseMargin_isGreatest, which do).  The top row's
% right label is the conjecture alpha_k(R) = 1, not a proved collapse; what is
% proved on that row is the bracket [1/k, 1] together with a design whose value
% is exactly 1 -- NOT that the infimum equals 1, which is the conjecture.
\begin{tikzpicture}[
  x=1cm, y=1cm,
  axis/.style={line width=0.5pt, black!70},
  minor/.style={line width=0.3pt, black!55},
  guide/.style={densely dotted, line width=0.3pt, black!35},
  thresh/.style={line width=0.5pt, black!55, dashed,
                 dash pattern=on 2.2pt off 1.6pt},
  span/.style={line width=0.9pt, black},
  runoff/.style={-{Latex[length=2mm,width=1.5mm]}, line width=0.9pt, black},
  fall/.style={-{Latex[length=2mm,width=1.5mm]}, line width=0.4pt, black!55},
  hair/.style={line width=0.4pt, black!45},
  every node/.style={font=\footnotesize}]

% ---- the guides, from the axis up through the rank rows ------------------
\foreach \gx in {2.00000, 2.66667, 5.61493, 6.11146, 6.76240, 7.51377}
  \draw[guide] (\gx,0.02) -- (\gx,3.38);

% ---- the threshold, at v = 1 --------------------------------------------
\draw[thresh] (8.00000,-0.14) -- (8.00000,3.52);
\node[anchor=south, font=\scriptsize] at (8.00000,3.50) {threshold};

% ======================= the value axis ==================================
\draw[axis] (-0.25,0) -- (9.75,0);
\foreach \mx in {0,0.8,1.6,2.4,3.2,4,4.8,5.6,6.4,7.2,8,8.8,9.6}
  \draw[minor] (\mx,0) -- (\mx,-0.07);

% named constants: a dot on the axis, a tick, and a label in one of two rows
\foreach \cx in {0.00000, 2.00000, 4.00000, 5.61493, 6.76240, 8.00000}{
  \draw[axis] (\cx,0) -- (\cx,-0.14);
  \fill (\cx,0) circle (1.3pt);}
\foreach \cx in {2.66667, 6.11146, 7.51377}{
  \draw[axis] (\cx,0) -- (\cx,-0.50);
  \fill (\cx,0) circle (1.3pt);}

\node[anchor=north] at (0.00000,-0.16) {$0$};                       % 0
\node[anchor=north] at (2.00000,-0.16) {$\tfrac14$};                % 1/4
\node[anchor=north] at (4.00000,-0.16) {$\tfrac12$};                % 1/2
\node[anchor=north] at (5.61493,-0.16) {$h_3$};                     % 3(1-cos 2pi/9)
\node[anchor=north] at (6.76240,-0.16) {$\alpha_{\mathrm{SIC}}$};   % 2-2/sqrt3
\node[anchor=north] at (8.00000,-0.16) {$1$};                       % 1
\node[anchor=north] at (2.66667,-0.52) {$\tfrac13$};                % 1/3
\node[anchor=north] at (6.11146,-0.52) {$3-\sqrt5$};                % 3-sqrt5
\node[anchor=north] at (7.51377,-0.52)
  {$\tfrac{10}9\alpha_{\mathrm{SIC}}$};                             % (10/9)(2-2/sqrt3)

% ======================= the six brackets ================================
% alpha_k(R), k >= 3 : [1/k, 1], right end attained, collapse conjectural
\draw[runoff] (2.66667,3.10) -- (0.92000,3.10);                     % 1/3, leftward
\draw[span]   (2.66667,3.10) -- (8.00000,3.10);                     % 1/3 .. 1
\draw[span]   (8.00000,3.025) -- (8.00000,3.175);
\fill (8.00000,3.10) circle (1.7pt);
\node[anchor=east] at (-0.30,3.10) {$\alpha_k(\R)$, $k\ge3$};
\node[anchor=west, font=\scriptsize] at (8.55,3.10) {$=1$ conjecturally};

% alpha_2(R) = 1
\draw[span] (7.84000,2.66) -- (8.16000,2.66);
\fill (8.00000,2.66) circle (1.7pt);
\node[anchor=east] at (-0.30,2.66) {$\alpha_2(\R)$};
\node[anchor=west, font=\scriptsize] at (8.55,2.66) {Theorem~\ref{thm:ranktwo}};

% alpha_k(C), k >= 4 : [1/k, (10/9) alpha_SIC]
\draw[runoff] (2.00000,2.10) -- (0.92000,2.10);                     % 1/4, leftward
\draw[span]   (2.00000,2.10) -- (7.51377,2.10);                     % 1/4 .. (10/9)aSIC
\draw[span]   (7.51377,2.025) -- (7.51377,2.175);
\node[anchor=east] at (-0.30,2.10) {$\alpha_k(\C)$, $k\ge4$};
\node[anchor=west, font=\scriptsize] at (8.55,2.10) {Theorem~\ref{thm:ladder}};

% alpha_3(C) : [1/3, h_3], the two superseded records chained above it
\draw[span] (2.66667,1.66) -- (5.61493,1.66);                       % 1/3 .. h_3
\draw[span] (2.66667,1.585) -- (2.66667,1.735);
\draw[span] (5.61493,1.585) -- (5.61493,1.735);
\fill (5.61493,1.66) circle (1.7pt);          % h_3 attained: Hesse design at (9,3)
\draw[fall] (7.51377,1.66) -- (5.72000,1.66);                       % (10/9)aSIC -> h_3
\foreach \ox in {6.11146, 7.51377}
  {\draw[line width=0.4pt, black!55, fill=white] (\ox,1.66) circle (1.5pt);}
\node[anchor=east] at (-0.30,1.66) {$\alpha_3(\C)$};
\node[anchor=west, font=\scriptsize] at (8.55,1.66) {open};

% alpha_2(C) = alpha_SIC
\draw[span] (6.60240,1.22) -- (6.92240,1.22);
\fill (6.76240,1.22) circle (1.7pt);
\node[anchor=east] at (-0.30,1.22) {$\alpha_2(\C)$};
\node[anchor=west, font=\scriptsize] at (8.55,1.22)
  {\S\ref{sec:cpadding} and \cite{Nesterenko2026complex}};

% alpha_1(F) = 1
\draw[span] (7.84000,0.78) -- (8.16000,0.78);
\fill (8.00000,0.78) circle (1.7pt);
\node[anchor=east] at (-0.30,0.78) {$\alpha_1(\FF)$};
\node[anchor=west, font=\scriptsize] at (8.55,0.78)
  {Remark~\ref{rem:sharpc}};

% ======================= the ledger ======================================
% Four columns at -2.15, -0.62, 4.16, 6.45, set from the measured cell widths
% (widest entry per column 0.98, 4.64, 2.14, 5.70cm); six rows 0.37 apart from -1.32.
\draw[hair] (-2.15,-1.05) -- (12.15,-1.05);
\begin{scope}[every node/.style={anchor=west, inner sep=1.5pt,
                                 font=\scriptsize}]
\node at (-2.15,-1.32) {$1/k$};
\node at (-0.62,-1.32) {floor at every rank, either field};
\node at ( 4.16,-1.32) {$\tfrac12,\ \tfrac13,\ \tfrac14,\ \ldots$};
\node at ( 6.45,-1.32) {Proposition~\ref{prop:rankfloor}};

\node at (-2.15,-1.69) {$h_3$};
\node at (-0.62,-1.69) {$3(1-\cos\tfrac{2\pi}9)$, least root of $h$};
\node at ( 4.16,-1.69) {$0.7018666706\ldots$};
\node at ( 6.45,-1.69)
  {Proposition~\ref{prop:hesse}, the Hesse design at $(9,3)$};

\node at (-2.15,-2.06) {$3-\sqrt5$};
\node at (-0.62,-2.06) {least root of $s^{2}-6s+4$};
\node at ( 4.16,-2.06) {$0.7639320225\ldots$};
\node at ( 6.45,-2.06)
  {Example~\ref{ex:trine}, two coplanar trines at $(6,3)$};

\node at (-2.15,-2.43) {$\alpha_{\mathrm{SIC}}$};
\node at (-0.62,-2.43) {$2-2/\sqrt3$, least root of $z^{2}-4z+\tfrac83$};
\node at ( 4.16,-2.43) {$0.8452994616\ldots$};
\node at ( 6.45,-2.43) {Lemma~\ref{lem:sicvalue}, the SIC at $(4,2)$};

\node at (-2.15,-2.80) {$\tfrac{10}9\alpha_{\mathrm{SIC}}$};
\node at (-0.62,-2.80) {$\tfrac{20}9-\tfrac{20\sqrt3}{27}$};
\node at ( 4.16,-2.80) {$0.9392216240\ldots$};
\node at ( 6.45,-2.80)
  {Theorem~\ref{thm:ladder}, the padded seed at $(k+2,k)$};

\node at (-2.15,-3.17) {$1$};
\node at (-0.62,-3.17) {the threshold, and $\alpha_1(\FF)$};
\node at ( 6.45,-3.17)
  {Theorem~\ref{thm:corankone}(iii)--(iv), attained at $(k+1,k)$};
\end{scope}

\draw[hair] (-2.15,-3.36) -- (12.15,-3.36);
\node[anchor=north west, align=left, text width=13.6cm,
      font=\scriptsize, inner sep=1.5pt] at (-2.15,-3.43)
  {$h(s)=8s^{3}-72s^{2}+162s-81$.  The roots of $h$ sum to $9$ and $h(2)=19>0$,
   so one root lies below $2$; with $h(3-\sqrt5)=14\sqrt5-27>0$ this gives
   $h_3<3-\sqrt5<2$.  Then $12<7\sqrt3$ gives
   $\alpha_{\mathrm{SIC}}>\tfrac56$, \ $5>(\tfrac{223}{100})^{2}$ gives
   $3-\sqrt5<\tfrac{77}{100}$, and $2500>2079$ closes
   $3-\sqrt5<\tfrac{77}{100}<\tfrac{25}{27}<\tfrac{10}9\alpha_{\mathrm{SIC}}$;
   \ and $363<400$ gives $\tfrac{10}9\alpha_{\mathrm{SIC}}<1$.};

\end{tikzpicture}
